\documentclass{ctexart}
\usepackage{template}
\usepackage{esint,extarrows,tikz}
\begin{document}\pagestyle{empty}
\begin{center}\Large
北京大学数学科学学院2022-23高等数学A1期末考试
\end{center}
\begin{problem}[1.(16\songti{分})]
回答下列问题。
        \begin{enumerate}[label=\tbf{(\arabic*)}]
            \item \textbf{(8\songti{分})}\ 证明:直线$l:\left\{\begin{array}{l}
                    x-2y+z=0\\
                    5x+2y-5z=-6
                \end{array}\right.$过点$(1,2,3)$,并把此一般方程化为标准方程.
            \item \textbf{(8\songti{分})}\ 求曲线$S:\left\{\begin{array}{l}
                    x=7t-14\\
                    y=4t^2\\
                    z=3t^3    
                \end{array}\right.$在参数$t=1$对应的点$P$处的法平面方程.
        \end{enumerate}
\end{problem}
\begin{solution}
\begin{enumerate}[label=\tbf{(\arabic*)},leftmargin=*]
\item 将$(1,2,3)$代入两平面方程,分别得到$0$和$-6$,故该点在直线上.两法向量的叉积为
\[(1,-2,1)\times(5,2,-5)=(8,10,12)\]
可取方向向量$(4,5,6)$,所以标准方程为
\[\dfrac{x-1}{4}=\dfrac{y-2}{5}=\dfrac{z-3}{6}\]
\item $P=(-7,4,3)$,曲线在$t=1$处的切向量为$(7,8,9)$.因此法平面为
\[7(x+7)+8(y-4)+9(z-3)=0\]
即$7x+8y+9z-10=0$.
\end{enumerate}
\end{solution}
\begin{problem}[2.(20\songti{分})]
回答下列问题.
        \begin{enumerate}[label=\tbf{(\arabic*)}]
            \item \textbf{(10\songti{分})}\ 设函数$z=\arctan\dfrac{(x-3)y+(x^2+x-1)y^2}{(x-2)y+(x-3)^2y^4}$,求$\left.\dfrac{\p z}{\p y}\right|_{(3,0)}$.
            \item \textbf{(10\songti{分})}\ 设函数$z=z(x,y)$由方程\[m\left(x+\dfrac zy\right)^n+n\left(y+\dfrac zx\right)^m=1\]确定,其中$m,n\in\N$.计算并化简\[x\dfrac{\p z}{\p x}+y\dfrac{\p z}{\p y}+xy\]
        \end{enumerate}
\end{problem}
\begin{solution}
\begin{enumerate}[label=\tbf{(\arabic*)},leftmargin=*]
\item 原题所写的分式在$(3,0)$处为$0/0$,函数在该点未定义,严格说来原题的偏导数不存在定义.
若按其在该点附近的连续延拓理解,对$y\neq0$约去公因子$y$,可得
\[\widetilde z(x,y)=\arctan\dfrac{x-3+(x^2+x-1)y}{x-2+(x-3)^2y^3}\]
延拓函数在$(3,0)$附近光滑,且$\widetilde z(3,y)=\arctan(11y)$,故
\[\widetilde z_y(3,0)=11\]
此数值对应连续延拓后的问题.
\item 在方程有定义且满足隐函数非退化条件的点上,令$w=z+xy$,则原方程变成
\[m(w/y)^n+n(w/x)^m=1\]
记$D=x\partial_x+y\partial_y$.对上式应用$D$,得到
\[mn\left(\dfrac{(w/y)^{n-1}}{y}+\dfrac{(w/x)^{m-1}}{x}\right)(Dw-w)=0\]
括号前的系数即原方程对$z$的偏导数,在非退化点不为零,于是$Dw=w$.
又$Dw=Dz+2xy$,故所求结果为
\[xz_x+yz_y+xy=z\]
若未假设非退化条件,应另核实隐函数的存在性与可微性.
\end{enumerate}
\end{solution}
\begin{problem}[3.(24\songti{分})]
下列极限是否存在?若存在,请求出其值;若不存在,请说明理由.
        \begin{enumerate}[label=\tbf{(\arabic*)}]
            \item \textbf{(8\songti{分})}\ \(\displaystyle\lim_{x\to0}\dfrac{\displaystyle\int_0^{x^3}\sin^32t\di t}{\displaystyle\int_0^{x^2}\tan t^5\di t}\).
            \item \textbf{(8\songti{分})}\ \(\displaystyle\lim_{(x,y)\to(0,0)}\dfrac{x^2y}{x^4+y^2}\).
            \item \textbf{(8\songti{分})}\ \(\displaystyle\lim_{n\to\infty}\sum_{k=1}^{n}\dfrac{n}{n^2+k^2}\).
        \end{enumerate}
\end{problem}
\begin{solution}
\begin{enumerate}[label=\tbf{(\arabic*)},leftmargin=*]
\item 当$t\to0$时,$\sin^3(2t)\sim8t^3$,$\tan(t^5)\sim t^5$.由积分的局部展开,
\[\int_0^{x^3}\sin^3(2t)\di t=2x^{12}+o(x^{12}),\qquad
\int_0^{x^2}\tan(t^5)\di t=\dfrac{x^{12}}6+o(x^{12})\]
因此极限为$12$.
\item 沿$y=0$趋近时函数值恒为$0$;沿$y=x^2$趋近时函数值恒为$1/2$.故二元极限不存在.
\item 将和式写成黎曼和,
\[\sum_{k=1}^n\dfrac{n}{n^2+k^2}
=\dfrac1n\sum_{k=1}^n\dfrac{1}{1+(k/n)^2}
\longrightarrow\int_0^1\dfrac{\di t}{1+t^2}=\dfrac\pi4\]
\end{enumerate}
\end{solution}
\begin{problem}[4.(24\songti{分})]
回答下列问题.
        \begin{enumerate}[label=\tbf{(\arabic*)}]
            \item \textbf{(8\songti{分})}\ 设$P_1(a_1,b_1,c_1),P_2(a_2,b_2,c_2)$是单位球面$S:x^2+y^2+z^2=1$上的两个不同的点,$O(0,0,0)$是坐标原点.求\[\left|\overrightarrow{OP_1}\times\overrightarrow{OP_2}\right|^2+\left(\overrightarrow{OP_1}\cdot\overrightarrow{OP_2}\right)\]
            \item \textbf{(8\songti{分})}\ 计算$\displaystyle\int_{-1}^1\left(\dfrac{\sin^2x}{1+\e^x}+\dfrac{\cos^2x}{1+\e^{-x}}\right)\dx$.
            \item \textbf{(8\songti{分})}\ 计算\(\displaystyle\int_{0}^{\frac{\pi}{4}}\ln\dfrac{\sin\left(x+\frac\pi4\right)}{\cos x}\dx\).
        \end{enumerate}
\end{problem}
\begin{solution}
\begin{enumerate}[label=\tbf{(\arabic*)},leftmargin=*]
\item 记$c=\overrightarrow{OP_1}\cdot\overrightarrow{OP_2}=a_1a_2+b_1b_2+c_1c_2$.
由拉格朗日恒等式及两向量长度均为$1$,有
\[|\overrightarrow{OP_1}\times\overrightarrow{OP_2}|^2=1-c^2\]
因此按现有题干,所求为$1-c^2+c$,一般不是常数.
若原意是点积的平方,则结果才恒为$1$;因未找到扫描原件,这里保留现有题干,不补写未经证实的平方.
\item 将积分中的第二项作代换$x=-t$,则$\cos^2(-t)=\cos^2t$,从而原积分等于
\[\int_{-1}^1\dfrac{\sin^2x+\cos^2x}{1+\e^x}\di x
=\int_{-1}^1\dfrac{\di x}{1+\e^x}\]
利用$\dfrac1{1+\e^x}+\dfrac1{1+\e^{-x}}=1$,得到积分值为$1$.
\item 对第一项令$u=x+\pi/4$,对第二项令$u=\pi/2-x$,则
\[\int_0^{\pi/4}\ln\sin(x+\pi/4)\di x
=\int_{\pi/4}^{\pi/2}\ln\sin u\di u
=\int_0^{\pi/4}\ln\cos x\di x\]
因此所求积分为$0$.
\end{enumerate}
\end{solution}
\begin{problem}[5.(8\songti{分})]
设$f(x)$在$(a,b)$上二阶可导,$f(a)=f(b)=0$,$f\left(\dfrac{a+b}{2}\right)>0$.试证明:存在$\xi\in(a,b)$使得$f''(\xi)<0$.
\end{problem}
\begin{solution}
现有题干只假设开区间内二阶可导,未假设在端点连续,按字面命题不成立.
例如令$f(x)=1$对$x\in(a,b)$成立,并单独指定$f(a)=f(b)=0$,则$f''\equiv0$.
若补充$f\in C[a,b]$,令$c=(a+b)/2$,由两次Lagrange中值定理,
存在$u\in(a,c)$和$v\in(c,b)$使
\[f'(u)=\dfrac{f(c)}{c-a}>0,\qquad f'(v)=-\dfrac{f(c)}{b-c}<0\]
再对$f'$在$[u,v]$上应用中值定理,得某$\xi\in(u,v)$满足
\[f''(\xi)=\dfrac{f'(v)-f'(u)}{v-u}<0\]
故补充端点连续条件后结论成立.
\end{solution}
\begin{problem}[6.(8\songti{分})]
设$f(x)$在$[0,2]$上有连续的导数,$f(0)=f(2)=0$,记$M=\displaystyle\max_{x\in[0,2]}\left\{|f(x)|\right\}$.试证明:
        \begin{enumerate}[label=\tbf{(\arabic*)}]
            \item \textbf{(4\songti{分})}\ 存在$\xi\in(0,2)$使得$\left|f'(\xi)\right|\geqslant M$.
            \item \textbf{(4\songti{分})}\ 若对于任意$x\in(0,2)$都有$\left|f'(x)\right|\leqslant M$,则$f(x)\equiv0$.
        \end{enumerate}
\end{problem}
\begin{solution}
\begin{enumerate}[label=\tbf{(\arabic*)},leftmargin=*]
\item 若$M=0$,则$f\equiv0$,结论显然.若$M>0$,取$c\in(0,2)$使$|f(c)|=M$.
对$[0,c]$和$[c,2]$分别应用中值定理,得到导数绝对值分别为$M/c$和$M/(2-c)$的两个点.
由于$\min(c,2-c)\leqslant1$,其中至少一个导数的绝对值不小于$M$.
\item 反设$M>0$,必要时将$f$替换成$-f$,取$c\in(0,2)$使$f(c)=M$.此时$f'(c)=0$.
若$c\leqslant1$,由$f'$连续且$|f'|\leqslant M$,在$c$附近有严格小于$M$的上界,所以
\[M=f(c)-f(0)=\int_0^c f'(x)\di x<Mc\leqslant M\]
矛盾.若$c\geqslant1$,同理有$M=-\int_c^2 f'(x)\di x<M(2-c)\leqslant M$,仍矛盾.
因此$M=0$,即$f\equiv0$.
\end{enumerate}
\end{solution}
\end{document}
